% ============================================================
%  Chapter 3 -- Convergent Sequences and Subsequences, 3.1-3.7
% ============================================================
\rchapter{3}{Numerical Sequences and Series}

\begin{roadmap}[Why this chapter exists]
Chapter 1 promised that limits would have somewhere to arrive; Chapter 2 built
the spaces they arrive in. This chapter finally defines the arrival ---
$p_n \to p$ --- and then spends its second half on the single most important
special case: series, which are sequences of partial sums.

\medskip
Two theorems are the chapter's spine. \textbf{3.11} says that in $\R^k$ a
sequence converges exactly when it is a Cauchy sequence --- convergence
becomes an \emph{internal} property, checkable without knowing the limit.
\textbf{3.14} says a monotonic sequence converges exactly when it is bounded.
Both are completeness in action, and nearly every convergence test for series
(3.24 through 3.44) is one of them wearing different clothes.
\end{roadmap}

\begin{roadmap}[How it is organized, and why in that order]
\textbf{Rudin's own account}, from the opening paragraph: the chapter ``will
deal primarily with sequences and series of complex numbers,'' but ``the basic
facts about convergence are just as easily explained in a more general
setting'' --- so the first three sections work in metric spaces, and the
theorems proved there are inherited free by $\R$, $\C$ and $\R^k$. The same
economy as 1.13(c) and Chapter 2.

\medskip
\centering
\begin{tikzpicture}[scale=0.9, transform shape, node distance=4mm]
\node[rmbox, text width=20mm] (a) {\textbf{3.1--3.20}\\sequences:\\limits, Cauchy,\\$\limsup$};
\node[rmbox, text width=20mm, right=of a] (b) {\textbf{3.21--3.32}\\series; the\\number $e$};
\node[rmbox, text width=20mm, right=of b] (c) {\textbf{3.33--3.46}\\tests; power\\series};
\node[rmbox, text width=20mm, right=of c] (d) {\textbf{3.47--3.55}\\products, re-\\arrangements};
\draw[rmarrow] (a) -- (b); \draw[rmarrow] (b) -- (c); \draw[rmarrow] (c) -- (d);
\end{tikzpicture}

\medskip
\raggedright
Three things to know in advance.

\smallskip
\textbf{The sequence sections are the load-bearing wall.} Everything about
series is a statement about the sequence of partial sums, so 3.21--3.55 consume
3.1--3.20 continuously. In particular $\limsup$ (3.16--3.17), the chapter's
hardest definition, exists mainly so that the root test (3.33) and the radius
of convergence (3.39) can be stated without exceptions.

\smallskip
\textbf{Completeness gets its name here.} Definition 3.12 finally introduces
the word \emph{complete} for ``every Cauchy sequence converges,'' discharging
the promise made in the comparison box at Definition 1.10. Two more of the
ECNU text's six equivalent completeness statements appear as theorems in this
chapter: monotone convergence is 3.14, and the Cauchy criterion is 3.11(c).

\smallskip
\textbf{The last section is a warning.} Rearranging a series can change its
sum (3.54) unless the convergence is absolute (3.55). The chapter is arranged
so that this lands last, as the reason the distinction of 3.45--3.46 was worth
drawing at all.
\end{roadmap}

\noindent
As the title indicates, this chapter will deal primarily with sequences and
series of complex numbers. The basic facts about convergence, however, are just
as easily explained in a more general setting. The first three sections will
therefore be concerned with sequences in euclidean spaces, or even in metric
spaces.

\rsection{Convergent Sequences}

\ritem{3.1}{Definition}
A sequence $(p_n)$ in a metric space $X$ is said to \emph{converge} if there is
a point $p \in X$ with the following property: For every $\varepsilon > 0$
there is an integer $N$ such that $n \ge N$ implies that
$d(p_n,p) < \varepsilon$. (Here $d$ denotes the distance in $X$.)

In this case, we also say that $(p_n)$ \emph{converges to} $p$, or that $p$ is
the \emph{limit} of $(p_n)$ [see Theorem 3.2(b)], and we write $p_n \to p$, or
\[ \lim_{n\to\infty} p_n = p. \]
If $(p_n)$ does not converge, it is said to \emph{diverge}.\kn{Read the
quantifiers in the stated order: $\varepsilon$ first, then $N$. The challenge
comes before the response, and $N$ may depend on $\varepsilon$ --- almost
always does. Every limit argument in this book is a game of respond-to-%
$\varepsilon$, and the bracketed reference to 3.2(b) is Rudin flagging that
``\emph{the} limit'' needs the uniqueness proved there before the phrase is
legitimate.}

It might be well to point out that our definition of ``convergent sequence''
depends not only on $(p_n)$ but also on $X$; for instance, the sequence
$(1/n)$ converges in $\R^1$ (to $0$), but fails to converge in the set of all
positive real numbers [with $d(x,y) = |x-y|$]. In cases of possible ambiguity,
we can be more precise and specify ``convergent in $X$'' rather than
``convergent.''\kw{The same relativity as ``open'' in 2.29, and the same cure:
name the ambient space. The sequence $(1/n)$ has not changed; what changed is
whether the space still contains the point it is crowding toward. Completeness
(3.12) is exactly the property of spaces in which this failure mode ---
converging toward a missing point --- cannot occur.}

We recall that the set of points $p_n$ $(n = 1,2,3,\dots)$ is the \emph{range}
of $(p_n)$. The range of a sequence may be a finite set, or it may be infinite.
The sequence $(p_n)$ is said to be \emph{bounded} if its range is bounded.

As examples, consider the following sequences of complex numbers (that is,
$X = \R^2$):
\begin{enumerate}[label=(\alph*), leftmargin=2.2em, itemsep=2pt, topsep=4pt]
  \item If $s_n = 1/n$, then $\lim_{n\to\infty} s_n = 0$; the range is
        infinite, and the sequence is bounded.
  \item If $s_n = n^2$, the sequence $(s_n)$ is unbounded, is divergent, and
        has infinite range.
  \item If $s_n = 1 + [(-1)^n/n]$, the sequence $(s_n)$ converges to $1$, is
        bounded, and has infinite range.
  \item If $s_n = i^n$, the sequence $(s_n)$ is divergent, is bounded, and has
        finite range.\kn{The four powers of $i$ recur forever, so the sequence
        visits $\{i,-1,-i,1\}$ in rotation. It diverges even though every point
        of its range is visited infinitely often --- each of the four points is
        a \emph{subsequential} limit (3.5), and 3.2(b) says a convergent
        sequence can have only one.}
  \item If $s_n = 1$ $(n = 1,2,3,\dots)$, then $(s_n)$ converges to $1$, is
        bounded, and has finite range.
\end{enumerate}

We now summarize some important properties of convergent sequences in metric
spaces.

\ritem{3.2}{Theorem}
\emph{Let $(p_n)$ be a sequence in a metric space $X$.}
\begin{enumerate}[label=(\alph*), leftmargin=2.2em, itemsep=1pt, topsep=3pt]
  \item \emph{$(p_n)$ converges to $p \in X$ if and only if every neighborhood
        of $p$ contains $p_n$ for all but finitely many $n$.}
  \item \emph{If $p \in X$, $p' \in X$, and if $(p_n)$ converges to $p$ and to
        $p'$, then $p' = p$.}
  \item \emph{If $(p_n)$ converges, then $(p_n)$ is bounded.}
  \item \emph{If $E \subset X$ and if $p$ is a limit point of $E$, then there
        is a sequence $(p_n)$ in $E$ such that $p = \lim_{n\to\infty} p_n$.}
\end{enumerate}

\begin{proof}
\begin{enumerate}[label=(\alph*), leftmargin=2.2em, itemsep=4pt, topsep=3pt]
  \item Suppose $p_n \to p$ and let $V$ be a neighborhood of $p$. For some
        $\varepsilon > 0$, the conditions $d(q,p) < \varepsilon$, $q \in X$
        imply $q \in V$. Corresponding to this $\varepsilon$, there exists $N$
        such that $n \ge N$ implies $d(p_n,p) < \varepsilon$. Thus $n \ge N$
        implies $p_n \in V$.

        Conversely, suppose every neighborhood of $p$ contains all but
        finitely many of the $p_n$. Fix $\varepsilon > 0$ and let $V$ be the
        set of all $q \in X$ such that $d(p,q) < \varepsilon$. By assumption,
        there exists $N$ (corresponding to this $V$) such that $p_n \in V$ if
        $n \ge N$. Thus $d(p_n,p) < \varepsilon$ if $n \ge N$; hence
        $p_n \to p$.
  \item Let $\varepsilon > 0$ be given. There exist integers $N$, $N'$ such
        that
        \begin{align*}
          n \ge N &\quad\text{implies}\quad d(p_n,p) < \frac{\varepsilon}{2},\\
          n \ge N' &\quad\text{implies}\quad d(p_n,p') < \frac{\varepsilon}{2}.
        \end{align*}
        Hence if $n \ge \max\{N,N'\}$, we have
        \[ d(p,p') \le d(p,p_n) + d(p_n,p') < \varepsilon. \]
        Since $\varepsilon$ was arbitrary, we conclude that
        $d(p,p') = 0$.\kn{``Since $\varepsilon$ was arbitrary'' is a standard
        closing move worth naming: a fixed nonnegative number smaller than
        every positive $\varepsilon$ must be $0$ (else take
        $\varepsilon = d(p,p')/2$). Rudin will use this silently from now on;
        Exercise 1.5.7 of Chapter 1's exercises --- Tao states it as a lemma
        --- is the same fact.}
  \item Suppose $p_n \to p$. Then there is an integer $N$ such that $n > N$
        implies $d(p_n,p) < 1$. Put
        \[ r = \max\{1,\, d(p_1,p),\, \dots,\, d(p_N,p)\}. \]
        Then $d(p_n,p) \le r$ for $n = 1,2,3,\dots$.\kn{The finite-maximum
        move of 2.20 and 2.24(c) again, from the other side: the tail is
        controlled by convergence, and only \emph{finitely many} early terms
        remain, so their distances have a maximum. Convergence controls tails;
        finiteness mops up beginnings.}
  \item For each positive integer $n$, there is a point $p_n \in E$ such that
        $d(p_n,p) < 1/n$. Given $\varepsilon > 0$, choose $N$ so that
        $N\varepsilon > 1$. If $n > N$, it follows that
        $d(p_n,p) < \varepsilon$. Hence $p_n \to p$.
\end{enumerate}
This completes the proof.
\end{proof}

\begin{center}
\begin{tikzpicture}[x=1mm, y=1mm]
  % epsilon tube picture for convergence
  \draw[axis] (-2,0) -- (96,0);
  \node[lbl, anchor=north] at (94,-2) {$n$};
  % limit line and tube
  \draw[thin] (0,14) -- (92,14);
  \node[mlbl, anchor=east] at (-2,14) {$p$};
  \draw[guide] (0,19) -- (92,19);
  \draw[guide] (0,9) -- (92,9);
  \node[mlbl, anchor=west] at (92.5,19) {$p+\varepsilon$};
  \node[mlbl, anchor=west] at (92.5,9) {$p-\varepsilon$};
  % terms: scattered early, inside tube after N
  \foreach \x/\y in {4/26, 9/4, 14/22, 19/7, 24/24, 29/10, 34/20}{
    \node[ptfill, minimum size=2.2pt] at (\x,\y) {};}
  \foreach \x/\y in {44/16.5, 50/12, 56/15.8, 62/12.9, 68/15.1, 74/13.4, 80/14.7, 86/13.8}{
    \node[ptfill, minimum size=2.2pt] at (\x,\y) {};}
  \draw[guide] (39,-5) -- (39,28);
  \node[lbl, anchor=north] at (39,-6) {$N$};
  \node[lbl, anchor=south] at (18,29) {finitely many strays allowed};
  \node[lbl, anchor=south] at (66,21.5) {all terms from $N$ on};
\end{tikzpicture}
\end{center}
\figcap{Convergence, drawn for a real sequence: every $\varepsilon$-tube about
$p$ traps the whole tail. Shrink $\varepsilon$ and $N$ moves right; that is the
dependence the definition quantifies.}

For sequences in $\R^k$ we can study the relation between convergence, on the
one hand, and the algebraic operations on the other. We first consider
sequences of complex numbers.

\ritem{3.3}{Theorem}
\emph{Suppose $(s_n)$, $(t_n)$ are complex sequences and}
\[ \lim_{n\to\infty} s_n = s, \qquad \lim_{n\to\infty} t_n = t. \]
\emph{Then}
\begin{enumerate}[label=(\alph*), leftmargin=2.2em, itemsep=2pt, topsep=3pt]
  \item $\displaystyle\lim_{n\to\infty} (s_n+t_n) = s+t$;
  \item $\displaystyle\lim_{n\to\infty} cs_n = cs$,\quad
        $\displaystyle\lim_{n\to\infty} (c+s_n) = c+s$, \emph{for any number
        $c$};
  \item $\displaystyle\lim_{n\to\infty} s_nt_n = st$;
  \item $\displaystyle\lim_{n\to\infty} (1/s_n) = 1/s$, \emph{provided
        $s_n \ne 0$ $(n = 1,2,3,\dots)$ and $s \ne 0$.}
\end{enumerate}

\begin{proof}
\begin{enumerate}[label=(\alph*), leftmargin=2.2em, itemsep=4pt, topsep=3pt]
  \item Given $\varepsilon > 0$, there exist integers $N_1$, $N_2$ such that
        \begin{align*}
          n \ge N_1 &\quad\text{implies}\quad |s_n-s| < \frac{\varepsilon}{2},\\
          n \ge N_2 &\quad\text{implies}\quad |t_n-t| < \frac{\varepsilon}{2}.
        \end{align*}
        If $N = \max\{N_1,N_2\}$, then $n \ge N$ implies
        \[ |(s_n+t_n)-(s+t)| \le |s_n-s| + |t_n-t| < \varepsilon. \]
        This proves (a). The proof of (b) is trivial.
  \item[(c)] We use the identity
        \begin{equation}
          s_nt_n - st = (s_n-s)(t_n-t) + s(t_n-t) + t(s_n-s). \tag{3.1}
        \end{equation}\kn{The identity is the whole proof, and it is chosen so
        that every term on the right involves at least one \emph{error}
        factor $s_n-s$ or $t_n-t$. Add and subtract $st_n$ (or $s_nt$) and you
        find it. The same add-and-subtract manufacture drives the proofs of
        4.4, 5.3(b) and the product rules everywhere.}
        Given $\varepsilon > 0$, there are integers $N_1$, $N_2$ such that
        \begin{align*}
          n \ge N_1 &\quad\text{implies}\quad |s_n-s| < \sqrt{\varepsilon},\\
          n \ge N_2 &\quad\text{implies}\quad |t_n-t| < \sqrt{\varepsilon}.
        \end{align*}
        If we take $N = \max\{N_1,N_2\}$, $n \ge N$ implies
        \[ |(s_n-s)(t_n-t)| < \varepsilon, \]
        so that
        \[ \lim_{n\to\infty} (s_n-s)(t_n-t) = 0. \]
        We now apply (a) and (b) to (3.1) and conclude that
        \[ \lim_{n\to\infty} (s_nt_n - st) = 0. \]
  \item[(d)] Choosing $m$ such that $|s_n-s| < |s|/2$ if $n \ge m$, we see
        that
        \[ |s_n| > \frac12|s| \qquad (n \ge m). \]\kn{The standard
        ``stay away from zero'' step: once $s_n$ is within $|s|/2$ of $s$, it
        cannot be nearer than $|s|/2$ to $0$ (reverse triangle inequality,
        Exercise 1.13). Without it, $1/s_n$ could blow up even with
        $s \ne 0$. Every quotient-limit proof in the book opens with this
        move.}
        Given $\varepsilon > 0$, there is an integer $N > m$ such that
        $n \ge N$ implies
        \[ |s_n-s| < \frac12|s|^2\varepsilon. \]
        Hence, for $n \ge N$,
        \[ \left|\frac{1}{s_n}-\frac{1}{s}\right|
           = \left|\frac{s_n-s}{s_ns}\right|
           < \frac{2}{|s|^2}|s_n-s| < \varepsilon. \]
\end{enumerate}
\end{proof}

\begin{center}
\begin{tikzpicture}[x=1mm, y=1mm]
  \fill[black!10] (38,0) rectangle (48,24);
  \fill[black!10] (0,24) rectangle (38,32);
  \fill[pattern=north east lines, pattern color=black!45]
    (38,24) rectangle (48,32);
  \draw[thin] (0,0) rectangle (48,32);
  \draw[thin] (0,0) rectangle (38,24);
  \node[mlbl] at (19,12) {$st$};
  \node[mlbl, rotate=90] at (43,12) {$t(s_n{-}s)$};
  \node[mlbl] at (19,28) {$s(t_n{-}t)$};
  \node[mlbl, anchor=west] at (52,28) {$(s_n{-}s)(t_n{-}t)$};
  \draw[thin] (49.5,28.6) -- (46,28.6);
  \node[mlbl, anchor=north] at (19,-1.6) {$s$};
  \node[mlbl, anchor=north] at (43,-1.6) {$s_n-s$};
  \node[mlbl, anchor=east] at (-1.6,12) {$t$};
  \node[mlbl, anchor=east] at (-1.6,28) {$t_n-t$};
\end{tikzpicture}
\end{center}
\figcap{Identity (3.1) as areas: the outer rectangle $s_nt_n$ minus the inner
$st$ is two thin strips plus one hatched corner --- and each of the three
pieces carries at least one error factor. The strips are bounded${}\times{}$%
small, the corner is small${}\times{}$small, which is why $\sqrt\varepsilon$
suffices for it. (Drawn for real $s_n > s > 0$, $t_n > t > 0$; the identity is
pure algebra and holds for complex sequences.)}

\begin{unpack}[How these proofs are found, not just checked]
All four parts follow one recipe, and it is the recipe for every
$\varepsilon$-argument in the book. Write down the quantity that must be small
--- $|(s_n+t_n)-(s+t)|$, say. Bound it by pieces you control, each going to
$0$. Then walk the proof backwards: decide what each piece must be smaller
than so the total is under $\varepsilon$, and let that dictate the $N$'s.

The visible fingerprints of working backwards are the odd-looking thresholds:
$\varepsilon/2$ in (a) because two errors will be added;
$\sqrt{\varepsilon}$ in (c) because two errors will be \emph{multiplied};
$\frac12|s|^2\varepsilon$ in (d) because the error will be divided by
$|s_ns| > \frac12|s|^2$. None of these is guessed --- each is the final
inequality, solved for the intermediate one. This is the same
reverse-engineering as the choice of $h$ in 1.21.
\end{unpack}

\ritem{3.4}{Theorem}
\begin{enumerate}[label=(\alph*), leftmargin=2.2em, itemsep=3pt, topsep=3pt]
  \item \emph{Suppose $\mathbf{x}_n \in \R^k$ $(n = 1,2,3,\dots)$ and}
        \[ \mathbf{x}_n = (\alpha_{1,n},\dots,\alpha_{k,n}). \]
        \emph{Then $(\mathbf{x}_n)$ converges to
        $\mathbf{x} = (\alpha_1,\dots,\alpha_k)$ if and only if}
        \begin{equation}
          \lim_{n\to\infty} \alpha_{j,n} = \alpha_j \qquad (1 \le j \le k).
          \tag{3.2}
        \end{equation}
  \item \emph{Suppose $(\mathbf{x}_n)$, $(\mathbf{y}_n)$ are sequences in
        $\R^k$, $(\beta_n)$ is a sequence of real numbers, and
        $\mathbf{x}_n \to \mathbf{x}$, $\mathbf{y}_n \to \mathbf{y}$,
        $\beta_n \to \beta$. Then}
        \[ \lim_{n\to\infty}(\mathbf{x}_n+\mathbf{y}_n) = \mathbf{x}+\mathbf{y},
           \qquad
           \lim_{n\to\infty} \mathbf{x}_n\cdot\mathbf{y}_n =
           \mathbf{x}\cdot\mathbf{y}, \qquad
           \lim_{n\to\infty} \beta_n\mathbf{x}_n = \beta\mathbf{x}. \]
\end{enumerate}

\begin{proof}
\begin{enumerate}[label=(\alph*), leftmargin=2.2em, itemsep=4pt, topsep=3pt]
  \item If $\mathbf{x}_n \to \mathbf{x}$, the inequalities
        \[ |\alpha_{j,n}-\alpha_j| \le |\mathbf{x}_n-\mathbf{x}|, \]
        which follow immediately from the definition of the norm in $\R^k$,
        show that (3.2) holds.

        Conversely, if (3.2) holds, then to each $\varepsilon > 0$, there
        corresponds an integer $N$ such that $n \ge N$ implies
        \[ |\alpha_{j,n}-\alpha_j| < \frac{\varepsilon}{\sqrt k}
           \qquad (1 \le j \le k). \]
        Hence $n \ge N$ implies
        \[ |\mathbf{x}_n-\mathbf{x}| =
           \left\{\sum_{j=1}^{k} |\alpha_{j,n}-\alpha_j|^2\right\}^{1/2}
           < \varepsilon, \]
        so that $\mathbf{x}_n \to \mathbf{x}$. This proves (a).

        Part (b) follows from (a) and Theorem 3.3.\kn{Part (a) is the
        coordinatewise reduction of 2.39 again: convergence in $\R^k$ is $k$
        convergences in $\R$, squeezed between the two norm inequalities
        $|\alpha_j| \le |\mathbf{x}| \le \sqrt{k}\max_j|\alpha_j|$. It is why
        the rest of the chapter can work with complex or real sequences and
        lose nothing.}
\end{enumerate}
\end{proof}

\rsection{Subsequences}

\begin{signpost}[Why subsequences, and why now]
A subsequence discards terms but keeps order, and the notion pays for itself
twice. First, it converts Chapter 2's compactness into a statement about
sequences: 3.6(a) is Bolzano--Weierstrass in sequence form, the tool that
manufactures limits out of mere boundedness. Second, subsequential limits are
the raw material for $\limsup$ and $\liminf$ (3.16), which the series tests
will need. Theorem 3.7 --- the set of subsequential limits is closed --- is
stocked for 3.17.
\end{signpost}

\ritem{3.5}{Definition}
Given a sequence $(p_n)$, consider a sequence $(n_k)$ of positive integers
such that $n_1 < n_2 < n_3 < \cdots$. Then the sequence $(p_{n_i})$ is called a
\emph{subsequence} of $(p_n)$. If $(p_{n_i})$ converges, its limit is called a
\emph{subsequential limit} of $(p_n)$.

It is clear that $(p_n)$ converges to $p$ if and only if every subsequence of
$(p_n)$ converges to $p$. We leave the details of the proof to the
reader.\kd{One direction is trivial; the other is Exercise-%
level but worth doing once: since $n_k \ge k$ (induction on the strictly
increasing indices), the tail of a subsequence is trapped by the tail of the
sequence. The statement is used constantly in the contrapositive: to show
$(p_n)$ diverges, exhibit two subsequences with different limits --- as with
$i^n$ in 3.1(d).}

\ritem{3.6}{Theorem}
\begin{enumerate}[label=(\alph*), leftmargin=2.2em, itemsep=2pt, topsep=3pt]
  \item \emph{If $(p_n)$ is a sequence in a compact metric space $X$, then some
        subsequence of $(p_n)$ converges to a point of $X$.}
  \item \emph{Every bounded sequence in $\R^k$ contains a convergent
        subsequence.}
\end{enumerate}

\begin{proof}
\begin{enumerate}[label=(\alph*), leftmargin=2.2em, itemsep=4pt, topsep=3pt]
  \item Let $E$ be the range of $(p_n)$. If $E$ is finite then there is a
        $p \in E$ and a sequence $(n_i)$ with $n_1 < n_2 < n_3 < \cdots$ such
        that
        \[ p_{n_1} = p_{n_2} = \cdots = p. \]
        The subsequence $(p_{n_i})$ so obtained converges evidently to
        $p$.\kn{The finite case is not a formality --- it is the pigeonhole:
        infinitely many indices land on finitely many values, so some value is
        hit infinitely often. Rudin splits by the \emph{range}, the
        set-versus-sequence distinction he flagged at 2.7.}

        If $E$ is infinite, Theorem 2.37 shows that $E$ has a limit point
        $p \in X$. Choose $n_1$ so that $d(p,p_{n_1}) < 1$. Having chosen
        $n_1,\dots,n_{i-1}$, we see from Theorem 2.20 that there is an integer
        $n_i > n_{i-1}$ such that $d(p,p_{n_i}) < 1/i$. Then $(p_{n_i})$
        converges to $p$.\kn{Theorem 2.20 --- every neighbourhood of a limit
        point contains \emph{infinitely many} points of $E$ --- is exactly what
        lets the index increase: within distance $1/i$ there are infinitely
        many terms, so one of them has index beyond $n_{i-1}$. The $q \ne p$
        clause of 2.18(b) has been earning interest since Chapter 2.}
  \item This follows from (a) since Theorem 2.41 implies that every bounded
        subset of $\R^k$ lies in a compact subset of $\R^k$.
\end{enumerate}
\end{proof}

\begin{deeper}[3.6(b) is Bolzano--Weierstrass in sequence form]
Compare 2.42: every bounded infinite \emph{set} has a limit point. Part (b) is
the same fact retold for sequences, and it is the workhorse version --- most of
analysis manufactures a convergent subsequence from boundedness and then works
with it. The ECNU text lists this (\cjk{致密性定理}{the ``compactness theorem''}, its
name for sequential compactness) as one of its six equivalent forms of
completeness, and Tao takes
exactly this property as his \emph{definition} of compactness (his 1.5.1).

Note what the proof consumed: 2.37, hence compactness of some $k$-cell, hence
--- through 2.40 and 2.38 --- the least-upper-bound property. Over $\Q$ the
statement fails: the sequence $1, 1.4, 1.41, \dots$ of Chapter 1 is bounded
with no convergent subsequence in $\Q$.
\end{deeper}

\ritem{3.7}{Theorem}
\emph{The subsequential limits of a sequence $(p_n)$ in a metric space $X$ form
a closed subset of $X$.}

\begin{proof}
Let $E^*$ be the set of all subsequential limits of $(p_n)$ and let $q$ be a
limit point of $E^*$. We have to show that $q \in E^*$.

Choose $n_1$ so that $p_{n_1} \ne q$. (If no such $n_1$ exists, then $E^*$ has
only one point, and there is nothing to prove.) Put $\delta = d(q,p_{n_1})$.
Suppose $n_1,\dots,n_{i-1}$ are chosen. Since $q$ is a limit point of $E^*$,
there is an $x \in E^*$ with $d(x,q) < 2^{-i}\delta$. Since $x \in E^*$, there
is an $n_i > n_{i-1}$ such that $d(x,p_{n_i}) < 2^{-i}\delta$. Thus
\[ d(q,p_{n_i}) \le 2^{1-i}\delta \]
for $i = 1,2,3,\dots$. This says that $(p_{n_i})$ converges to $q$. Hence
$q \in E^*$.\kn{A two-step chase: $q$ is approximated by subsequential limits
$x$, and each $x$ is approximated by actual terms $p_{n_i}$ --- so terms of the
sequence approach $q$ through the middlemen, at the price of a doubled error
$2\cdot2^{-i}\delta$. The geometric budget $2^{-i}\delta$ is chosen so the
doubling still goes to zero. This diagonal-through-approximators pattern
returns at 7.23.}
\end{proof}
